\documentclass[a4paper]{article}
\usepackage{amsfonts}
\usepackage{amsmath}
\usepackage{amssymb}
\usepackage{graphicx}     

\begin{document}
  
\noindent \large Auteur: Michael Zielinski \\
\noindent \large Studierichting: Informatica\\
\noindent \large Oefeningenreeks 6A nr. 13\\

\medskip

\normalsize

Noem de drie oorspronkelijke getallen: \\

$ a,\ b,\ c $ \\

Omdat ze een meetkundige rij vormen, geldt: \\

$ b^2 = ac $ \\

Als men bij het middelste getal $2$ optelt, wordt de rij rekenkundig: \\

$ a,\ b+2,\ c $ \\

Dus: \\

$ 2(b+2)=a+c $ \\

$ 2b+4=a+c $ \\

Daarna telt men opnieuw $2$ bij het middelste getal op en $20$ bij het laatste getal: \\

$ a,\ b+4,\ c+20 $ \\

Deze rij is opnieuw meetkundig, dus: \\

$ (b+4)^2=a(c+20) $ \\

We krijgen het stelsel: \\

$ 
\left\{
\begin{array}{l}
b^2=ac \\
2b+4=a+c \\
(b+4)^2=a(c+20)
\end{array}
\right.
$ \\

Uitwerking geeft: \\

$ (a,b,c)=(4,8,16) $ \\

of \\

$ (a,b,c)=\left(\dfrac{4}{9},-\dfrac{8}{9},\dfrac{16}{9}\right) $ \\

De oorspronkelijke rij is dus: \\

$ \boxed{4,\ 8,\ 16} $ \\

of \\

$ \boxed{\dfrac{4}{9},\ -\dfrac{8}{9},\ \dfrac{16}{9}} $ \\

\begin{thebibliography}{999}
%\bibitem{a} Referentie
\end{thebibliography}

\end{document}